\begin{minipage}{0.78\linewidth}
    On dispose d'un produit cosmétique conçu pour le traitement de l'acné et de
    la rosacée en application topique. L'étiquette du produit porte l'indication
    \textbf{<<Acide Azélaïque 10\%>>} (Fig.~\ref{fig:1}). L'indication $10\%$
    désigne le pourcentage massique de l'acide Azélaïque dans le produit.

    À fin de vérifier l'indication précédente, on prépare une solution aqueuse
    $\mathrm{(S)}$ de concentration molaire $C$ et de volume $V=\qty{100}{\mL}$
    en dissolvant la masse ${m_{\text{(produit)}}=\qty{1.88}{\g}}$ du produit
    cosmétique dans l'eau distillée.

    On réalise le dosage du volume $V_{A}=\qty{10}{\mL}$ de la solution
    $\mathrm{(S)}$ par une solution aqueuse $\mathrm{(S_{B})}$ d'hydroxyde de
    sodium $\ce{Na^+_{(aq)} + HO^-_{(aq)}}$ de concentration molaire
    ${C_{B}=\qty{1e-2}{\mol\per\L}}$ en utilisant le montage de la
    Fig.~\ref{fig:2}.  La courbe de la Fig.~\ref{fig:3} représente les
    variations du $\mathrm{pH}$ du mélange en fonction du volume $V_{B}$ de la
    solution $\mathrm{(S_{B})}$ ajoutée.
\end{minipage}
\hfill
\begin{minipage}{0.20\linewidth}
    \begin{figure}[H]
        \centering
        \includegraphics[width=\textwidth]{2bac_svt_national_2024_rattrapage_chim_1}
        \caption{}
        \label{fig:2bac_svt_national_2024_rattrapage_chim_1}
    \end{figure}
\end{minipage}

\begin{minipage}{0.42\linewidth}
    \begin{figure}[H]
        \centering
        \includegraphics[width=\textwidth]{2bac_svt_national_2024_rattrapage_chim_2}
        \caption{}
        \label{fig:2bac_svt_national_2024_rattrapage_chim_2}
    \end{figure}
\end{minipage}
\hfill
\begin{minipage}{0.42\linewidth}
    \begin{figure}[H]
        \centering
        \includegraphics[width=\textwidth]{2bac_svt_national_2024_rattrapage_chim_3}
        \caption{}
        \label{fig:2bac_svt_national_2024_rattrapage_chim_3}
    \end{figure}
\end{minipage}

\begin{enumerate}
    \item Attribuer au numéros~(1),~(2) et~(3) du montage de la Fig.~\ref{fig:2}
          les noms correspondants.
    \item Écrire l'équation de la réaction du dosage supposée totale.
    \item Déterminer les coordonnées $(V_{B\mathrm{E}}~;~\mathrm{pH_{E}})$ du
          point d'équivalence.
    \item Calculer la concentration $C$ de la solution $\mathrm{(S)}$.
    \item Vérifier l'indication <<Acide Azélaïque $10\%$>> indiquée sur
          l'étiquette du produit cosmétique.
\end{enumerate}

